\documentclass[10pt]{beamer}
\usepackage{tikz}
\usepackage{tikz-3dplot}
\usetikzlibrary{positioning}
\usepackage{pgfplots}
\usepackage{booktabs}
\usepackage{multirow}
\usepackage{array}
\usepackage{siunitx}
\usepackage{makecell}
\usepackage{CJKutf8}
\usepackage{hyperref}
\usepackage{bookmark}

\pgfplotsset{compat=newest}
\setbeamertemplate{navigation symbols}{}

% Add page number to the lower right corner
\setbeamertemplate{footline}{%
  \hfill\usebeamercolor[fg]{page number in head/foot}%
  \usebeamerfont{page number in head/foot}%
  \insertframenumber\,/\,\inserttotalframenumber\hspace*{1ex}\vskip2pt%
}
\setbeamercolor{page number in head/foot}{fg=gray}
\setbeamerfont{page number in head/foot}{size=\small}

\title{Web Scraping \& APIs in Python}
\author{Tzu-You Lin}
\date{February 17, 2025}

% Automatically add a section title slide at the beginning of each section.
\AtBeginSection[]{
  \begin{frame}
    \vfill
    \centering
    \begin{beamercolorbox}[sep=8pt,center,shadow=true,rounded=true]{title}
      \usebeamerfont{title}\insertsectionhead\par%
    \end{beamercolorbox}
    \vfill
  \end{frame}
}

\begin{document}

% Title Page
\begin{frame}[fragile]
    \titlepage
\end{frame}

% Outline
\begin{frame}[fragile]{Outline}
    \tableofcontents
\end{frame}

%%%%%%%%%%%%%%%%%%%%%%%%%%%%%%%%%%%%%%%%%%%%%%%%%%%%%%%%%%%%%%%%%%%%%%
\section{Introduction}
\begin{frame}[fragile]{What is Web Scraping?}
    \begin{itemize}
        \item Automated extraction of data from websites.
        \item Involves fetching web pages and parsing the content.
        \item Tools include \texttt{Requests}, \texttt{BeautifulSoup}, \texttt{Scrapy}, and \texttt{Selenium}.
    \end{itemize}
\end{frame}

\begin{frame}[fragile]{What are APIs?}
    \begin{itemize}
        \item Application Programming Interfaces allow communication between software applications.
        \item Provide structured data (e.g., JSON, XML) from web services.
        \item Commonly used to access platforms like GitHub, Twitter, and many others.
    \end{itemize}
\end{frame}

\begin{frame}[fragile]{Web Scraping vs. APIs}
    \begin{itemize}
        \item \textbf{Web Scraping:}
        \begin{itemize}
            \item Extracts data from raw HTML.
            \item Often unstructured and subject to website layout changes.
        \end{itemize}
        \item \textbf{APIs:}
        \begin{itemize}
            \item Provides structured and well-documented data.
            \item Typically more stable and legally supported.
        \end{itemize}
        \item Choice depends on data availability and project requirements.
    \end{itemize}
\end{frame}

%%%%%%%%%%%%%%%%%%%%%%%%%%%%%%%%%%%%%%%%%%%%%%%%%%%%%%%%%%%%%%%%%%%%%%
\section{Tools and Techniques}
\begin{frame}[fragile]{Web Scraping Tools in Python}
    \begin{itemize}
        \item \textbf{Requests:} For making HTTP requests.
        \item \textbf{BeautifulSoup:} For parsing and navigating HTML/XML.
        \item \textbf{Scrapy:} A robust framework for large-scale scraping projects.
        \item \textbf{Selenium:} For automating browsers and scraping dynamic content.
    \end{itemize}
\end{frame}

\begin{frame}[fragile]{APIs in Python}
    \begin{itemize}
        \item \textbf{Requests:} Simplifies HTTP API calls.
        \item \textbf{JSON:} For parsing and handling JSON data.
        \item \textbf{Flask/Django:} To create your own API endpoints.
        \item \textbf{Async Libraries:} Such as \texttt{httpx} and \texttt{aiohttp} for asynchronous API calls.
    \end{itemize}
\end{frame}

\begin{frame}[fragile]{Making HTTP Requests}
    \begin{verbatim}
import requests

url = 'https://example.com/data'
response = requests.get(url)

if response.status_code == 200:
    data = response.json()
    print(data)
else:
    print("Error:", response.status_code)
    \end{verbatim}
\end{frame}

\begin{frame}[fragile]{Parsing HTML with BeautifulSoup}
    \begin{verbatim}
from bs4 import BeautifulSoup
import requests

url = 'https://www.example.com'
response = requests.get(url)
soup = BeautifulSoup(response.text, 'html.parser')

# Extract page title
title = soup.find('title').get_text()
print("Page Title:", title)
    \end{verbatim}
\end{frame}

\begin{frame}[fragile]{Error Handling in Scraping}
    \begin{itemize}
        \item Always check HTTP status codes.
        \item Use \texttt{try-except} blocks to catch exceptions.
        \item Implement logging to help diagnose issues.
    \end{itemize}
\end{frame}

\begin{frame}[fragile]{Respecting Website Policies}
    \begin{itemize}
        \item Review the site's \texttt{robots.txt} file to understand crawling rules.
        \item Use rate limiting to avoid overloading servers.
        \item Follow the website's terms of service.
    \end{itemize}
\end{frame}

%%%%%%%%%%%%%%%%%%%%%%%%%%%%%%%%%%%%%%%%%%%%%%%%%%%%%%%%%%%%%%%%%%%%%%
\section{APIs in Action}
\begin{frame}[fragile]{Working with APIs: GitHub Example}
    \begin{verbatim}
import requests

url = 'https://api.github.com/repos/python/cpython'
response = requests.get(url)
data = response.json()

print("Repository:", data['name'])
print("Stars:", data['stargazers_count'])
    \end{verbatim}
\end{frame}

\begin{frame}[fragile]{Working with APIs: Random Stuff Example}
    \begin{verbatim}
import requests

url = 'https://api.random.org/json-rpc/4/invoke'
payload = {
    "jsonrpc": "2.0",
    "method": "generateIntegers",
    "params": {
        "apiKey": "YOUR-API-KEY",
        "n": 5,
        "min": 1,
        "max": 100
    },
    "id": 1
}
response = requests.post(url, json=payload)

if response.status_code == 200:
    data = response.json()
    numbers = data['result']['random']['data']
    print("Random numbers:", numbers)
    \end{verbatim}
\end{frame}

\begin{frame}[fragile]{API Response Structure}
    \begin{itemize}
        \item The response contains:
        \begin{itemize}
            \item Random numbers generated
            \item Request metadata
            \item Usage information
            \item API version details
        \end{itemize}
        \item Example response:
        \begin{verbatim}
{
    "jsonrpc": "2.0",
    "result": {
        "random": {
            "data": [42, 23, 87, 15, 91],
            "completionTime": "2024-02-17 10:15:32Z"
        },
        "bitsUsed": 24,
        "requestsLeft": 999
    },
    "id": 1
}
        \end{verbatim}
    \end{itemize}
\end{frame}

\begin{frame}[fragile]{Data Extraction and Transformation}
    \begin{itemize}
        \item Clean and structure raw data.
        \item Use Pandas for data manipulation.
        \item Convert data into formats like CSV, JSON, or databases.
    \end{itemize}
\end{frame}

\begin{frame}[fragile]{Storing and Processing Data}
    \begin{itemize}
        \item Save data to CSV files using Pandas.
        \item Consider using databases (e.g., SQLite, PostgreSQL) for larger datasets.
        \item Visualize data with libraries like Matplotlib or Seaborn.
    \end{itemize}
\end{frame}

\begin{frame}[fragile]{Real-World Applications}
    \begin{itemize}
        \item Market research and competitive analysis.
        \item Social media sentiment analysis.
        \item Price tracking and trend monitoring.
        \item Academic research and report generation.
    \end{itemize}
\end{frame}

%%%%%%%%%%%%%%%%%%%%%%%%%%%%%%%%%%%%%%%%%%%%%%%%%%%%%%%%%%%%%%%%%%%%%%
\section{Advanced Topics and Best Practices}
\begin{frame}[fragile]{Challenges in Web Scraping}
    \begin{itemize}
        \item Handling dynamic, JavaScript-driven pages.
        \item Navigating anti-scraping measures (e.g., CAPTCHAs).
        \item Keeping up with frequent website changes.
    \end{itemize}
\end{frame}

\begin{frame}[fragile]{Using Selenium for Dynamic Content}
    \begin{itemize}
        \item Automates browsers to render JavaScript.
        \item Captures content that isn't loaded on the initial HTML.
        \item Example code:
    \end{itemize}
    \begin{verbatim}
from selenium import webdriver

driver = webdriver.Chrome()
driver.get("https://www.example.com")
html = driver.page_source
driver.quit()
    \end{verbatim}
\end{frame}

\begin{frame}[fragile]{API Best Practices}
    \begin{itemize}
        \item Use proper authentication methods (API keys, OAuth).
        \item Respect API rate limits and guidelines.
        \item Implement error handling and retries.
        \item Always consult and follow the API documentation.
    \end{itemize}
\end{frame}

\begin{frame}[fragile]{Ethics and Legal Considerations}
    \begin{itemize}
        \item Understand the legal boundaries of web scraping.
        \item Respect user privacy and intellectual property.
        \item When available, prefer using official APIs over scraping.
    \end{itemize}
\end{frame}

\begin{frame}[fragile]{Summary and Conclusion}
    \begin{itemize}
        \item Web scraping and APIs are vital tools for data extraction.
        \item Select the right approach based on data structure and legality.
        \item Follow best practices to ensure ethical and efficient operations.
        \item Experiment with advanced tools for complex scraping tasks.
    \end{itemize}
\end{frame}

\begin{frame}[fragile]{Further Resources}
    \begin{itemize}
        \item \textbf{Book:} \emph{Web Scraping with Python} by Ryan Mitchell.
        \item \textbf{Course Supplementary:} This course has supplementary materials on the course website.
        \item \textbf{Online Tutorials:} Real Python, Scrapy and BeautifulSoup documentation.
        \item \textbf{Communities:} Stack Overflow, GitHub projects.
        \item \textbf{APIs:} Official documentation for GitHub, Twitter, etc.
    \end{itemize}
\end{frame}

\end{document}
